@@M0@@

The proof, immediate from the functorial definition of flag schemes, is left to the reader.

\sgaStableTarget{sga6.e06.lem.4-3}\textbf{Lemma 4.3.} \uline{Let} @@M1@@ \uline{be a commutative ring with unit}, @@M2@@ \uline{an integer} @@M3@@, @@M4@@ \uline{a family of elements of} @@M5@@; \uline{let} @@M6@@ \uline{be the} @@M7@@-\uline{algebra generated by generators} @@M8@@, @@M9@@, \uline{subject to the relations}:
@@M10@@
\uline{where the} @@M11@@ \uline{denote the elementary symmetric functions}. \uline{Let} @@M12@@ \uline{be integers} @@M13@@ \uline{with sum} @@M14@@, \uline{and for} @@M15@@,
@@M16@@
\uline{with}
@@M17@@
\uline{Put}:
@@M18@@
\uline{and let} @@M19@@ \uline{be the subalgebra of} @@M20@@ \uline{generated by the} @@M21@@. \uline{Then}:
\begin{enumerate}[label=\textup{\roman*)}]
\item \uline{the elements of the form}:
@@M22@@
\uline{where} @@M23@@, and @@M24@@, \uline{form a basis of} @@M25@@ \uline{considered as a} @@M26@@-\uline{module}.
\item \uline{the generators} @@M27@@ \uline{of the} @@M28@@-\uline{algebra} @@M29@@ \uline{satisfy the} @@M30@@ \uline{isobaric relations deduced from the relation}:
@@M31@@
